\title{Controlling Text-to-Image Diffusion by Orthogonal Finetuning}

\begin{abstract}
Large-scale text-to-image diffusion models demonstrate remarkable ability to produce photorealistic images from textual descriptions. An important unresolved issue is how to effectively steer or control these powerful models to accomplish various downstream tasks. To address this, we propose a principled finetuning approach called Orthogonal Finetuning~(OFT) for adapting text-to-image diffusion models to specific downstream applications. In contrast to existing techniques, OFT can theoretically maintain hyperspherical energy, which quantifies the pairwise neuron relationships on the unit hypersphere. We observe that preserving this property is essential for maintaining the semantic generation capabilities of text-to-image diffusion models. To enhance finetuning stability, we additionally introduce Constrained Orthogonal Finetuning (COFT), which enforces an extra radius constraint on the hypersphere. In particular, we focus on two key finetuning scenarios for text-to-image models: subject-driven generation, where the objective is to generate subject-specific images based on a few images of the subject along with a text prompt, and controllable generation, where the goal is to enable the model to incorporate extra control signals. Our experimental results demonstrate that the OFT framework surpasses existing methods in terms of generation quality and convergence speed.
\end{abstract}

\section{Introduction}

Recent text-to-image diffusion models~\cite{saharia2022photorealistic,ramesh2022hierarchical,rombach2022high} have shown outstanding performance in text-guided generation of high-quality images. Despite these achievements, text prompts can still be ambiguous and insufficient for providing precise, fine-grained control over the generated images. To address this limitation, this work focuses on two categories of text-to-image generation tasks:

\begin{itemize}[leftmargin=*,nosep]

    \item \textbf{Subject-driven generation}~\cite{ruiz2023dreambooth}: Given only a few images of a particular subject, the task involves generating images of the same subject placed in different contexts, guided by a text prompt.
    \item \textbf{Controllable generation}~\cite{zhang2023adding,mou2023t2i}: Given an additional control input (e.g., canny edges, segmentation maps), the objective is to produce images that comply with both the control signal and the textual description.
\end{itemize}

Both tasks fundamentally revolve around how to efficiently finetune text-to-image diffusion models without compromising the generative capabilities acquired during pretraining. We outline the key requirements for an effective finetuning approach as: (1) \emph{training efficiency}: minimizing the number of trainable parameters and training epochs, and (2) \emph{generalizability preservation}: maintaining high-fidelity and diverse generative outputs. To achieve this, finetuning is commonly performed either by slightly adjusting neuron weights with a low learning rate (\emph{e.g.}, \cite{ruiz2023dreambooth}) or by incorporating a small additional module with re-parameterized neuron weights (\emph{e.g.}, \cite{hulora2022,zhang2023adding}). Despite their straightforwardness, neither approach guarantees retention of the generative performance established during pretraining. Moreover, there is a lack of systematic insight into designing effective finetuning strategies or determining appropriate hyperparameters like the training epoch count. A major challenge is the absence of a metric to quantify how well the pretrained generative ability is preserved. Current finetuning techniques implicitly assume that a smaller Euclidean distance between the finetuned and pretrained models corresponds to better preservation of pretrained capabilities. For this reason, finetuning methods typically employ very low learning rates. While this assumption can sometimes be valid, we contend that Euclidean distance alone fails to adequately reflect the extent of semantic preservation. Consequently, employing a more structurally informed metric to characterize differences between the finetuned and pretrained models can significantly enhance the maintenance of pretrained performance as well as the stability of finetuning.

Motivated by the empirical finding that hyperspherical similarity effectively encodes semantic information~\cite{liu2017deep,liu2018decoupled,chen2020angular}, we employ hyperspherical energy~\cite{liu2018learning} to represent the pairwise relational structure among neurons. Hyperspherical energy is defined as the aggregate of hyperspherical similarity (\emph{e.g.}, cosine similarity) between every pair of neurons within the same layer, reflecting the degree of neuron uniformity on the unit hypersphere~\cite{liu2021learning}. We propose that an optimal finetuned model should exhibit minimal deviation in hyperspherical energy relative to the pretrained model. A straightforward approach would be to introduce a regularizer that maintains the hyperspherical energy during finetuning; however, this does not guarantee effective minimization of the hyperspherical energy difference. Consequently, we leverage a key invariance property of hyperspherical energy—pairwise hyperspherical similarity is theoretically maintained when an identical orthogonal transformation is applied to all neurons. Inspired by this invariance, we introduce Orthogonal Finetuning~(OFT), a method that adapts large text-to-image diffusion models to downstream tasks without altering their hyperspherical energy. The core concept involves learning a layer-shared orthogonal transformation for the neurons, preserving their pairwise angular relationships. OFT can also be interpreted as modifying the canonical coordinate system for neurons within the same layer. Recognizing that a smaller Euclidean distance between the finetuned and pretrained models correlates with better retention of pretrained performance, we further develop a variant called Constrained Orthogonal Finetuning~(COFT), which restricts the finetuned model to lie within a hypersphere of fixed radius centered around the pretrained neurons.

The reasoning behind why orthogonal transformation is effective for finetuning neurons partly derives from the 2D Fourier transform, which decomposes an image into magnitude and phase spectra. The phase spectrum, representing the angular relationship between the input and basis, retains most of the semantic content. For instance, an image's phase spectrum combined with a random magnitude spectrum can still reconstruct the original image without losing its semantic meaning. This observation implies that altering neuron directions is essential for semantically modifying the generated image, which is the objective in both subject-driven and controllable generation. However, allowing unrestricted changes in neuron directions will inevitably impair the generative performance obtained from pretraining. To limit this freedom, we propose preserving the angle between any pair of neurons, grounded in the hypothesis that the angles among neurons are vital for encoding the neural network's knowledge. Based on this insight, it is intuitive to learn a layer-shared orthogonal transformation for neurons within each layer, ensuring the hyperspherical energy remains constant.

Additionally, we take inspiration from orthogonal over-parameterized training~\cite{liu2021orthogonal}, which trains classification neural networks from scratch by applying orthogonal transformations to a randomly initialized network. This approach is motivated by the fact that a randomly initialized neural network has, in expectation, a provably small hyperspherical energy, and the aim of \cite{liu2021orthogonal} is to maintain this small hyperspherical energy throughout training (as low hyperspherical energy correlates with better generalization in classification~\cite{liu2018learning,lin2020regularizing}). The work of \cite{liu2021orthogonal} demonstrates that orthogonal transformations provide enough flexibility to train neural networks that generalize well for classification tasks. By contrast, our focus is on finetuning text-to-image diffusion models to enhance controllability and improve downstream generative performance. We highlight two key distinctions between OFT and \cite{liu2021orthogonal}. First, whereas \cite{liu2021orthogonal} aims to minimize hyperspherical energy, OFT strives to preserve the pretrained model's hyperspherical energy to avoid disrupting its inherent structure during finetuning. Minimizing hyperspherical energy in diffusion model finetuning might damage the original semantic relationships. Second, OFT seeks to minimize deviation from the pretrained model, resulting in a constrained variant, whereas \cite{liu2021orthogonal} imposes no such restrictions. The essence of effective finetuning is balancing flexibility and stability, and we contend that our OFT framework successfully attains this balance. Our contributions are summarized as follows:

\begin{itemize}[leftmargin=*,nosep]

    \item We introduce a new finetuning technique called Orthogonal Finetuning (OFT) aimed at enhancing the controllability of text-to-image diffusion models. To further bolster stability, we present a constrained version that restricts the angular deviation from the pretrained model.
    \item Unlike existing finetuning approaches, OFT carries out model finetuning while theoretically guaranteeing the preservation of hyperspherical energy, which we empirically observe to be a crucial indicator of the pretrained model’s ability to maintain generative semantics.
    \item We apply OFT to two scenarios: subject-driven generation and controllable generation. Through extensive experiments, we show notable improvements over previous methods in terms of generation quality, convergence speed, and finetuning robustness. Additionally, OFT exhibits better sample efficiency, converging effectively with substantially fewer training images and epochs.
    \item For the controllable generation task, we propose a novel control consistency metric to quantitatively assess controllability. The central concept involves estimating the control signal from the generated output and comparing it to the original control input. This metric further confirms the superior controllability achieved by OFT.
\end{itemize}

\section{Related Work}

\textbf{Text-to-image diffusion models}. Significant advancements~\cite{saharia2022photorealistic,ramesh2022hierarchical,rombach2022high,gu2022vector,nichol2022glide} in text-to-image synthesis have been driven by the rapid progress in diffusion-based generative models~\cite{ho2020denoising,song2021score,song2021denoising,dhariwal2021diffusion} and vision-language representation learning~\cite{su2020vlbert,lu2019vilbert,radford2021learning,wang2021simvlm,li2022blip,li2023blip,alayrac2022flamingo,schuhmann2022laion}. GLIDE~\cite{nichol2022glide} and Imagen~\cite{saharia2022photorealistic} train diffusion models directly in pixel space. GLIDE jointly trains a text encoder with a diffusion prior using paired text-image datasets, whereas Imagen employs a frozen, pretrained text encoder. Stable Diffusion~\cite{rombach2022high} and DALL-E2~\cite{ramesh2022hierarchical} operate by training diffusion models within latent spaces. Stable Diffusion leverages VQ-GAN~\cite{esser2021taming} to learn a visual codebook latent space, while DALL-E2 uses CLIP~\cite{radford2021learning} to create a joint latent space embedding images and text. Besides diffusion models, other approaches like generative adversarial networks~\cite{reed2016generative,zhang2017stackgan,xu2018attngan,li2019controllable} and autoregressive models~\cite{ramesh2021zero,ding2021cogview,wu2022nuwa,yuscaling} have also been explored for text-to-image generation. OFT is fundamentally a model-agnostic finetuning technique applicable to any text-to-image diffusion model.

\textbf{Subject-driven generation}. To avoid altering the subject, \cite{avrahami2022blended,nichol2022glide} use user-provided masks as additional conditioning. Inversion techniques~\cite{choi2021ilvr,dhariwal2021diffusion,ramesh2022hierarchical,gal2022image} allow modifying the context while preserving the subject. \cite{hertz2022prompt} enables both local and global editing without requiring input masks. However, these methods often fail to adequately retain identity-specific details of the subject. Pivotal Tuning~\cite{roich2022pivotal} finetunes the generator around an initial inverted latent code with extra regularization to maintain identity. Similarly, \cite{nitzan2022mystyle} learns a personalized generative prior from a set of a person's facial images. \cite{casanova2021instance} generates various versions of an instance but may lose fine instance-specific details. Using a customized token and a limited set of subject images, DreamBooth~\cite{ruiz2023dreambooth} finetunes text-to-image diffusion models with both a reconstruction loss and a class-specific prior preservation loss. OFT builds upon the DreamBooth framework but diverges by employing orthogonal transformations for finetuning instead of simple finetuning with a reduced learning rate.

\textbf{Controllable generation}. Image-to-image translation can be regarded as a type of controllable generation, with prior approaches predominantly utilizing conditional generative adversarial networks~\cite{isola2017image,zhu2017toward,wang2018high,park2019semantic,choi2018stargan}. Diffusion models have also been employed for image-to-image translation tasks~\cite{saharia2022palette,voynov2022sketch,wang2022pretraining}. More recently, ControlNet~\cite{zhang2023adding} introduces a method to steer a pretrained diffusion model by finetuning and adapting it to incorporate extra control signals, achieving remarkable results in controllable generation. A concurrent and related approach, T2I-Adapter~\cite{mou2023t2i}, similarly finetunes a pretrained diffusion model to enhance controllability over generated images. Following the same task framework as \cite{zhang2023adding,mou2023t2i}, we utilize OFT to finetune pretrained diffusion models, consistently attaining improved controllability with fewer training samples and fewer parameters updated during finetuning. Importantly, OFT does not add any extra computational burden during inference.

\textbf{Model finetuning}. The practice of finetuning large pretrained models for downstream applications has become increasingly widespread~\cite{devlin2018bert,brown2020language,he2022masked}. Among finetuning strategies, adaptation techniques (such as \cite{hulora2022,houlsby2019parameter,pfeiffer2020adapterfusion}) have been extensively explored in natural language processing. LoRA~\cite{hulora2022} is the most closely related method to OFT, positing a low-rank structure for the additive weight modifications during finetuning. In contrast, OFT employs a layer-shared orthogonal transformation to update neuron weights multiplicatively, and it is theoretically shown to preserve the pairwise angles between neurons within the same layer, leading to enhanced stability.

\section{Orthogonal Finetuning}

\subsection{Why Does Orthogonal Transformation Make Sense?}

We begin by examining why orthogonal transformation is advantageous for finetuning text-to-image diffusion models. This question is divided into two components: (1) why it is beneficial to finetune the angles (i.e., directions) of neurons, and (2) why orthogonal transformation is chosen as the means to finetune these angles.

For the initial question, we take inspiration from the empirical findings in \cite{liu2018decoupled,chen2020angular} which indicate that angular differences in features effectively represent the semantic gap. SphereNet~\cite{liu2017deep} demonstrates that training a neural network with all neurons normalized onto a unit hypersphere provides similar capacity and even enhances generalization, suggesting that the directions of neurons can capture the most critical information from the data. To further illustrate the significance of neuron angles, we conduct a simple experiment shown in Figure~\ref{fig:toy_ae}, where we train a conventional convolutional autoencoder on a set of flower images. During training, the feature map is generated using the standard inner product ($z$ represents the output element of the convolution kernel $\bm{w}$, and $\bm{x}$ denotes the input in the sliding window). At testing, we compare three methods for producing the feature map: (a) the inner product applied during training, (b) magnitude information, and (c) angular information. The results in Figure~\ref{fig:toy_ae} reveal that the angular information of neurons can nearly perfectly reconstruct the input images, whereas the magnitude of neurons holds no meaningful information. It is important to note that we do not apply the cosine activation (c) during training, with training relying solely on the inner product. These findings suggest that neuron angles (directions) play a primary role in encoding the semantic content of input images. Consequently, adjusting neuron directions is likely to be a more effective way to alter the semantic information of images.

Regarding the second question, the most straightforward approach to fine-tune neuron directions is to simultaneously rotate or reflect all neurons within the same layer, inherently involving an orthogonal transformation. Although it could be more flexible to apply angular transformations that rotate different neurons by distinct angles, we find that orthogonal transformation strikes a balance between flexibility and regularity. Furthermore, \cite{liu2021orthogonal} demonstrates that orthogonal transformation possesses sufficient capacity for learning in neural networks. To substantiate our claim, we conduct an experiment illustrating the effective regularization induced by the orthogonality constraint. We carry out a controllable generation experiment based on the ControlNet~\cite{zhang2023adding} setup, with results displayed in Figure~\ref{fig:ortho}. Our observations indicate that the standard OFT performs consistently and achieves precise control after training (epoch 20). In contrast, OFT without the orthogonality constraint fails to generate any realistic images and does not produce any control effect. This experiment underscores the critical role of the orthogonality constraint in OFT.

\subsection{General Framework}

The standard finetuning approach generally employs gradient descent with a low learning rate to modify a model (or selected layers within a model). The use of a small learning rate implicitly promotes minimal deviation from the pretrained model, and conventional finetuning effectively attempts to update the model while implicitly minimizing $\thickmuskip=2mu \medmuskip=2mu \| \bm{M} - \bm{M}^0\|$, where $\bm{M}$ represents the finetuned model parameters and $\bm{M}^0$ denotes the pretrained parameters. This implicit restriction is intuitively reasonable but might still offer excessive flexibility when finetuning large models. To tackle this, LoRA adds an explicit low-rank constraint on the weight update, namely $\thickmuskip=2mu \medmuskip=2mu \text{rank}(\bm{M} - \bm{M}^0)=r'$, where $r'$ is chosen to be a small integer. In contrast, OFT imposes a constraint based on pairwise neuron similarity: $\thickmuskip=2mu \medmuskip=2mu \| \text{HE}(\bm{M})-\text{HE}(\bm{M}^0) \|=0$, where $\text{HE}(\cdot)$ denotes the hyperspherical energy of a weight matrix. To illustrate, consider a fully connected layer $\thickmuskip=2mu \medmuskip=2mu \bm{W}=\{\bm{w}_1,\cdots,\bm{w}_n\}\in\mathbb{R}^{d\times n}$ with neurons $\bm{w}_i\in\mathbb{R}^{d}$ ($\bm{W}^0$ is the pretrained weight matrix). The output vector $\thickmuskip=2mu \medmuskip=2mu \bm{z}\in\mathbb{R}^n$ is computed via $\thickmuskip=2mu \medmuskip=2mu \bm{z}=\bm{W}^\top\bm{x}$, where $\thickmuskip=2mu \medmuskip=2mu \bm{x}\in\mathbb{R}^d$ is the input vector. OFT can be understood as minimizing the difference in hyperspherical energy between the finetuned and pretrained models:

\begin{equation}\label{eq:oft_principle}
\footnotesize
\min_{\bm{W}} \left\| \text{HE}(\bm{W})-\text{HE}(\bm{W}^0)\right\|~~\Leftrightarrow~~\min_{\bm{W}} \bigg{\|} \sum_{i\neq j} \|\hat{\bm{w}}_i-\hat{\bm{w}}_j\|^{-1}-\sum_{i\neq j} \|\hat{\bm{w}}^0_i-\hat{\bm{w}}^0_j\|^{-1}\bigg{\|}
\end{equation}

where $\thickmuskip=2mu \medmuskip=2mu \hat{\bm{w}}_i:=\bm{w}_i/\|\bm{w}_i\|$ denotes the $i$-th neuron normalized to unit length, and the hyperspherical energy of a layer $\bm{W}$ is defined as $\thickmuskip=2mu \medmuskip=2mu \text{HE}(\bm{W}):= \sum_{i\neq j} \|\hat{\bm{w}}_i-\hat{\bm{w}}_j\|^{-1}$. It is straightforward to see that the minimum achievable value for Eq.~\eqref{eq:oft_principle} is zero. This minimum is attained whenever $\bm{W}$ differs from $\bm{W}^0$ solely by an orthogonal transformation, i.e., $\thickmuskip=2mu \medmuskip=2mu \bm{W}=\bm{R}\bm{W}^0$ where $\thickmuskip=2mu \medmuskip=2mu \bm{R}\in\mathbb{R}^{d\times d}$ is an orthogonal matrix (with determinant $1$ indicating a rotation and $-1$ indicating a reflection). This concept underpins OFT, which posits that finetuning a neural network can be accomplished by learning layer-shared orthogonal matrices that transform neurons within each layer. Formally, OFT aims to optimize the orthogonal

\begin{equation}\label{eq:oft_general}
    \footnotesize
    \bm{z}=\bm{W}^\top\bm{x}=(\bm{R}\cdot\bm{W}^0)^\top\bm{x},~~~\text{s.t.}~\bm{R}^\top\bm{R}=\bm{R}\bm{R}^\top=\bm{I}
\end{equation}

Here, $\bm{W}$ represents the OFT-finetuned weight matrix, and $\bm{I}$ is the identity matrix. OFT is depicted in Figure~\ref{fig:block}. Analogous to the zero initialization approach used in LoRA, it is necessary to ensure that OFT fine-tunes the pretrained model starting exactly from $\bm{W}^0$. To accomplish this, we initialize the orthogonal matrix $\bm{R}$ as an identity matrix, ensuring the fine-tuned model begins with the pretrained weights. To maintain the orthogonality of $\bm{R}$, differential orthogonalization techniques as presented in \cite{liu2021orthogonal,lezcano2019cheap} can be applied. We will elaborate on how to enforce orthogonality in a computationally efficient manner.

\subsection{Efficient Orthogonal Parameterization}\label{sect:eff_ortho}

Although standard orthogonalization methods like the Gram-Schmidt process are differentiable, they tend to be computationally intensive in practice~\cite{liu2021orthogonal}. To improve efficiency, we employ the Cayley parameterization to construct the orthogonal matrix. Concretely, the orthogonal matrix is formed as $\thickmuskip=2mu \medmuskip=2mu \bm{R}=(\bm{I}+\bm{Q})(\bm{I}-\bm{Q})^{-1}$, where $\bm{Q}$ is a skew-symmetric matrix satisfying $\thickmuskip=2mu \medmuskip=2mu \bm{Q}=-\bm{Q}^\top$. This approach offers greater efficiency at a minor cost: the Cayley parameterization only produces orthogonal matrices with determinant $1$, meaning it belongs to the special orthogonal group. Fortunately, this constraint does not negatively impact performance in practice. Despite using the Cayley transform for parameterization, $\bm{R}$ can remain parameter-inefficient when the dimension $d$ is large. To mitigate this, we propose to model $\bm{R}$ as a block-diagonal matrix comprising $r$ blocks, yielding the following structure:

\begin{equation}
    \footnotesize
    \bm{R}=\text{diag}(\bm{R}_1,\bm{R}_2,\cdots,\bm{R}_r)=\begin{bmatrix}
    \bm{R}_1\in \text{O}(\frac{d}{r}) &  &\\
    &\ddots & \\
    & & \bm{R}_r\in \text{O}(\frac{d}{r})
    \end{bmatrix}\in \text{O}(d)
\end{equation}

where $\text{O}(d)$ represents the orthogonal group in dimension $d$, and $\thickmuskip=2mu \medmuskip=2mu \bm{R}\in\mathbb{R}^{d\times d}$ and $\thickmuskip=2mu \medmuskip=2mu \bm{R}_i\in\mathbb{R}^{d/r\times d/r},\forall i$ are orthogonal matrices. When $\thickmuskip=2mu \medmuskip=2mu r=1$, the block-diagonal orthogonal matrix simplifies to a conventional unconstrained orthogonal matrix. For an orthogonal matrix of size $\thickmuskip=2mu \medmuskip=2mu d\times d$, the parameter count is $\thickmuskip=2mu \medmuskip=2mu d(d-1)/2$, corresponding to a computational complexity of $\mathcal{O}(d^2)$. In contrast, for an $r$-block diagonal orthogonal matrix, the number of parameters reduces to $\thickmuskip=2mu \medmuskip=2mu d(d/r-1)/2$, which leads to a complexity of $\thickmuskip=2mu \medmuskip=2mu \mathcal{O}(d^2/r)$. Additionally, parameters can be further decreased by sharing blocks, i.e., setting $\thickmuskip=2mu \medmuskip=2mu \bm{R}_i=\bm{R}_j$ for all $i \neq j$, reducing parameter complexity to $\mathcal{O}(d^2/r^2)$. Despite these parameter efficiency improvements, the resulting matrix $\bm{R}$ remains orthogonal, preserving hyperspherical energy without compromise.

We compare the parameter efficiency of OFT and LoRA. For LoRA with a low-rank parameter $r'$, the trainable parameter count is $\thickmuskip=2mu \medmuskip=2mu r'(d+n)$. Considering both $r$ and $r'$ as functions of the neuron dimension $d$ (e.g., $\thickmuskip=2mu \medmuskip=2mu r=r'=\alpha d$ for some constant $\thickmuskip=2mu \medmuskip=2mu 0<\alpha\leq 1$), LoRA’s parameter complexity is $\mathcal{O}(d^2 + dn)$, whereas OFT’s complexity is $\mathcal{O}(d)$. To illustrate this difference concretely, consider a weight matrix of size $\thickmuskip=2mu \medmuskip=2mu 128\times128$: LoRA has $2,048$ trainable parameters with $\thickmuskip=2mu \medmuskip=2mu r' = 8$, while OFT requires only $960$ parameters with $\thickmuskip=2mu \medmuskip=2mu r = 8$ (without block sharing).

\subsection{Constrained Orthogonal Finetuning}

The flexibility of the original OFT can be further controlled by restricting the finetuned model to remain close to the pretrained model. Specifically, COFT performs the forward pass as follows:

\begin{equation}\label{eq:coft_general}
    \footnotesize
    \bm{z} = \bm{W}^\top \bm{x} = (\bm{R} \cdot \bm{W}^0)^\top \bm{x}, \quad \text{s.t.} \quad \bm{R}^\top \bm{R} = \bm{R} \bm{R}^\top = \bm{I}, \quad \| \bm{R} - \bm{I} \| \leq \epsilon
\end{equation}

which imposes both an orthogonality constraint and an $\epsilon$-distance constraint to the identity matrix. The orthogonality can be enforced using the Cayley parameterization described in Section~\ref{sect:eff_ortho}. However, incorporating the $\epsilon$-deviation constraint into the Cayley-parameterized orthogonal matrix is challenging. To better understand the Cayley transform, we approximate $\thickmuskip=2mu \medmuskip=2mu \bm{R} = (\bm{I} + \bm{Q})(\bm{I} - \bm{Q})^{-1}$ via the Neumann series as $\thickmuskip=2mu \medmuskip=2mu \

(series converges in the operator norm). Consequently, we can incorporate the constraint $\thickmuskip=2mu \medmuskip=2mu\|\bm{R}-\bm{I}\|\leq\epsilon$ within the Cayley transform itself, leading to an equivalent constraint $\thickmuskip=2mu \medmuskip=2mu \|\bm{Q}-\bm{0}\|\leq\epsilon'$, where $\bm{0}$ represents the zero matrix and $\epsilon'$ is a different error hyperparameter from $\epsilon$. This updated constraint on the matrix $\bm{Q}$ can be straightforwardly enforced using projected gradient descent. To initialize the orthogonal matrix $\bm{R}$ as the identity, we set $\bm{Q}$ initially to the zero matrix. COFT can thus be interpreted as integrating two explicit constraints: minimizing the hyperspherical energy difference and restricting deviation from the pretrained model. While the latter is often implicitly incorporated in existing finetuning approaches, COFT explicitly enforces it. Although OFT performs excellently, we note that COFT offers superior finetuning stability compared to OFT due to this explicit deviation limitation. Figure~\ref{fig:eps_coft} illustrates how varying $\epsilon$ impacts COFT's performance. It shows that $\epsilon$ regulates finetuning flexibility: larger $\epsilon$ causes the COFT-finetuned model to resemble OFT-finetuned models more closely, whereas smaller $\epsilon$ makes the COFT-finetuned model behave increasingly like the pretrained text-to-image diffusion model.

\subsection{Re-scaled Orthogonal Finetuning}

We introduce a straightforward extension to the original OFT by incorporating a learnable magnitude scaling factor for each neuron. This is motivated by the observation that rescaling neurons does not affect the hyperspherical energy since magnitudes are normalized out. Specifically, the forward pass is defined as $\thickmuskip=2mu \medmuskip=2mu\bm{z}=(\bm{R}\bm{W}^0\bm{D})^\top\bm{x}$\footnote{\textbf{Errata}: In the NeurIPS camera ready version, the forward pass of re-scaled OFT was incorrectly stated as $\thickmuskip=2mu \medmuskip=2mu\bm{z}=(\bm{D}\bm{R}\bm{W}^0)^\top\bm{x}$. The original implementation remains correct, so results in Appendix~\ref{app:rescaled} are unaffected.}, where $\thickmuskip=2mu \medmuskip=2mu \bm{D}=\text{diag}(s_1,\cdots,s_n)\in\mathbb{R}^{n\times n}$ is a learnable diagonal matrix whose diagonal elements $s_1,\cdots,s_n$ are all positive. Unlike the original OFT forward pass in Eq.~\eqref{eq:oft_general} where only $\bm{R}$ is trainable, here both the diagonal matrix $\bm{D}$ and the orthogonal matrix $\bm{R}$ are learnable. The re-scaled OFT enhances the flexibility of OFT with a negligible increase in parameters. We retain the original OFT in experiments to demonstrate the efficacy of orthogonal transformation alone, but generally, we find re-scaled OFT performs better (refer to Appendix~\ref{app:rescaled}).

\section{Intriguing Insights and Discussions}

\textbf{OFT is architecture-agnostic}. In principle, OFT can be applied to any neural network architecture. For Transformers, LoRA is usually applied to attention weights~\cite{hulora2022}. To ensure fair comparison with LoRA, we apply OFT only to finetune attention weights in our experiments. Beyond fully connected layers, OFT is also well suited to finetuning convolutional layers due to the block-diagonal structure of $\bm{R}$, which has meaningful interpretations in convolutional layers (unlike LoRA). When using as many blocks as input channels, each block transforms a distinct neuron channel, analogous to learning depth-wise convolution kernels~\cite{chollet2017xception}. When all blocks in $\bm{R}$ are shared, OFT applies an orthogonal transformation shared across channels. We provide a preliminary study on finetuning convolutional layers with OFT in Appendix~\ref{app:convolution}.

\textbf{Relation to LoRA}. LoRA introduces a low-rank matrix to ensure that the pretrained weight matrix's information does not change significantly. In contrast, OFT enforces that the transformation applied to the pretrained weight matrix is orthogonal (full-rank), thereby preserving the pretrained information. The forward pass of OFT can be expressed as $\thickmuskip=2mu \medmuskip=2mu \bm{z} = (\bm{R}\bm{W}^0)^\top \bm{x} = (\bm{W}^0 + (\bm{R} - \bm{I}) \bm{W}^0)^\top \bm{x}$, where $\thickmuskip=2mu \medmuskip=2mu (\bm{R} - \bm{I}) \bm{W}^0$ functions similarly to LoRA's low-rank weight adjustment. Given that $\bm{W}^0$ is generally full-rank, OFT also performs a low-rank weight modification when $\thickmuskip=2mu \medmuskip=2mu \bm{R} - \bm{I}$ has low rank. Analogous to LoRA's rank parameter $r'$, OFT employs a diagonal block parameter $r$ to reduce the count of trainable parameters. Notably, LoRA and OFT embody two different strategies for parameter efficiency: LoRA leverages low-rank structure to decrease trainable parameters, while OFT utilizes sparsity structure (i.e., block-diagonal orthogonality).

\textbf{Reasons for OFT's faster convergence}. Firstly, as demonstrated in Figure~\ref{fig:toy_ae}, the most effective way to alter semantic representations is by changing neuron directions, which is precisely the design goal of OFT. Secondly, OFT can be interpreted as finetuning neurons constrained to a smooth hyperspherical manifold, which offers a more favorable optimization landscape. This observation is also supported by empirical results in \cite{liu2021orthogonal}.

\textbf{Why we do not minimize hyperspherical energy}. Unlike \cite{liu2021orthogonal}, our approach does not aim to minimize hyperspherical energy. In classification tasks, it is preferable to have neurons without redundancy. Minimizing hyperspherical energy corresponds to neurons being uniformly distributed on the hypersphere, which is not an appropriate objective for finetuning since it may disrupt the pretrained knowledge.

\textbf{Balancing flexibility and regularity in finetuning}. We identify a fundamental trade-off between flexibility and regularity. While standard finetuning offers the greatest flexibility, it tends to result in poor stability and can easily cause model collapse. Surprisingly simple, OFT strikes an effective balance between flexibility and regularity by maintaining the pairwise neuron angles. The block-diagonal parameterization can be interpreted as a stronger form of regularization applied to the orthogonal matrix.

\textbf{No extra inference cost}. Unlike ControlNet, our OFT framework does not add any inference overhead to the finetuned model. During inference, we can directly multiply the learned orthogonal matrix $\bm{R}$ with the pretrained weight matrix $\bm{W}^0$, producing an equivalent weight matrix $\thickmuskip=2mu \medmuskip=2mu \bm{W} = \bm{R}\bm{W}^0$. Therefore, the inference speed remains identical to that of the pretrained model.

\section{Experiments and Results}

\textbf{General settings}. In our experiments, we utilize Stable Diffusion v1.5~\cite{rombach2022high} as the pretrained text-to-image model. For fairness, we randomly select generated images from each method. For subject-driven generation, we largely follow DreamBooth~\cite{ruiz2023dreambooth}. For controllable generation, we mostly adhere to the protocols in ControlNet~\cite{zhang2023adding} and T2I-Adapter~\cite{mou2023t2i}. To enable a fair comparison with LoRA, we apply OFT or COFT only to the same layers where LoRA is implemented. Additional results and experimental details are provided in Appendix~\ref{app:settings}.

\subsection{Subject-driven Generation}

\textbf{Settings}. We compare against DreamBooth~\cite{ruiz2023dreambooth} and LoRA~\cite{hulora2022} baselines. All methods use the same loss function as in DreamBooth. For DreamBooth and LoRA, we generally follow their original implementations and apply the best hyperparameter configurations. Further results are included in Appendices~\ref{app:settings}, \ref{app:coft_oft}, \ref{app:more_qualitative}, and \ref{app:failure}.

\textbf{Stability and convergence of finetuning}. We begin by assessing the finetuning stability and convergence speed for DreamBooth, LoRA, OFT, and COFT. The results are presented in Figure~\ref{fig:teaser} and Figure~\ref{fig:db_convergence}. It is evident that both COFT and OFT can finetune the diffusion model in a stable manner. After 400 iterations, DreamBooth and OFT variants both attain good control, whereas LoRA fails to maintain the subject identity. At 2000 iterations, DreamBooth starts producing distorted images, and LoRA is unable to generate a yellow shirt, instead producing yellow fur. By contrast, OFT and COFT continue to provide stable and consistent control over the generated outputs. These findings confirm the rapid yet stable convergence of the OFT framework in subject-driven generation. We emphasize that the insensitivity to the number of finetuning iterations is crucial, as it effectively reduces the difficulty of tuning iteration counts for various subjects. For OFT and COFT, one can simply choose a relatively large iteration count without precise tuning. For COFT, with an appropriate $\epsilon$, setting both the learning rate and iteration number becomes straightforward.

\textbf{Quantitative evaluation}. Following \cite{ruiz2023dreambooth}, we perform a quantitative comparison measuring subject fidelity (DINO~\cite{caron2021emerging}, CLIP-I~\cite{radford2021learning}), text prompt fidelity (CLIP-T~\cite{radford2021learning}), and sample diversity (LPIPS~\cite{zhang2018unreasonable}). CLIP-I computes the average pairwise cosine similarity of CLIP embeddings between generated and real images. DINO is similar to CLIP-I, but uses ViT S/16 DINO embeddings instead. CLIP-T measures the average cosine similarity between CLIP embeddings of the text prompt and the generated images. We also calculate the average LPIPS cosine similarity between generated images of the same subject under the same text prompt. Table~\ref{table:dreambooth_final} demonstrates that both COFT and OFT outperform DreamBooth and LoRA substantially on the DINO and CLIP-I metrics, while achieving slightly better or comparable results in prompt fidelity and diversity. To minimize randomness, each method is finetuned repeatedly on the same pretrained model using 30 different random seeds. These results indicate that the OFT framework not only attains superior convergence and stability but also consistently delivers improved final performance.

\textbf{Qualitative comparison}. To gain a clearer insight into the advantages of OFT, we present several randomly selected examples of subject-driven generation in Figure~\ref{fig:dreambooth_final}. For an unbiased comparison, all examples are generated from the same finetuned model for each method, ensuring that no individual text prompt is separately optimized for its final output. For each technique, we choose the model that attains the highest validation CLIP metrics. From the outcomes shown in Figure~\ref{fig:dreambooth_final}, it is evident that both OFT and COFT maintain excellent semantic subject consistency, whereas LoRA frequently fails to retain the subject’s identity (e.g., LoRA completely loses the subject identity in the bowl example). Additionally, OFT and COFT demonstrate far more precise control via text prompts, whereas DreamBooth, although it preserves subject identity, often does not generate images aligned with the text prompt (e.g., the first row in the bowl example). This qualitative comparison highlights that our OFT framework simultaneously achieves superior controllability and subject preservation. Furthermore, OFT is robust to the number of iterations and performs well even with a high iteration count, unlike DreamBooth and LoRA. Additional qualitative examples are provided in Appendix~\ref{app:more_qualitative}. We also carry out a human evaluation in Appendix~\ref{app:human} which further confirms the superiority of OFT.

\subsection{Controllable Generation}

\textbf{Settings}. We adopt ControlNet~\cite{zhang2023adding}, T2I-Adapter~\cite{mou2023t2i}, and LoRA~\cite{hulora2022} as baseline methods. In the main paper, we focus on three challenging controllable generation tasks: converting Canny edge maps to images (C2I) on the COCO dataset~\cite{lin2014microsoft}, transforming segmentation maps to images (S2I) on the ADE20K dataset~\cite{zhou2017scene}, and mapping landmarks to face images (L2F) on the CelebA-HQ dataset~\cite{karras2018progressive,xia2021tedigan}. Each method finetunes Stable Diffusion (SD) v1.5 on these datasets for 20 epochs. Further results are available in Appendices~\ref{app:more_qualitative}, \ref{app:more_control_task}, and \ref{app:failure}.

\textbf{Convergence}. We assess the convergence rate of ControlNet, T2I-Adapter, LoRA, and COFT on the L2F task through both quantitative and qualitative analyses. For the quantitative assessment, we calculate the mean $\ell_2$ distance between the control face landmarks and the predicted face landmarks. In Figure~\ref{fig:face_convergence}, we display the face landmark error for models fine-tuned over varying numbers of epochs. It is evident that both COFT and OFT exhibit considerably faster convergence. Specifically, LoRA requires 20 epochs to reach the performance level that our OFT framework attains by the 8th epoch. We highlight that OFT and COFT utilize a comparable number of trainable parameters to LoRA (which is substantially fewer than ControlNet), yet they converge more efficiently than existing methods. Furthermore, the rapid convergence of OFT is corroborated by the results shown in Figure~\ref{fig:teaser}. The right-hand example in Figure~\ref{fig:teaser} illustrates that OFT is notably more data-efficient than ControlNet and LoRA, as it converges effectively using only 5\% of the ADE20K dataset. For qualitative evaluation, we concentrate on comparing OFT, COFT, and ControlNet, since ControlNet attains landmark errors closest to ours. The results in Figure~\ref{fig:face_convergence_qual} demonstrate that both OFT and COFT converge steadily, with the generated face pose progressively aligning with the control landmarks. In contrast to this smooth and stable convergence, ControlNet exhibits a sudden onset of controllability after the 8th epoch, which aligns with the abrupt convergence behavior reported in \cite{zhang2023adding}. This stable convergence pattern in our OFT framework makes it much more user-friendly in practice, as its training dynamics are significantly smoother and more predictable, thereby facilitating the selection of optimal OFT hyperparameters.

\textbf{Quantitative comparison}. We propose a control consistency metric to assess the effectiveness of controllable generation. The core concept involves extracting the control signal from the generated image and then comparing it against the original input control signal. For the C2I task, we measure IoU and F1 score. For the S2I task, we calculate mean IoU, mean accuracy, and overall accuracy. For the L2F task, we determine the mean $\ell_2$ distance between the control landmarks and the predicted landmarks. Additional details about these consistency metrics can be found in Appendix~\ref{app:settings}. For all the methods compared, we utilize the optimal hyperparameter configurations. The results shown in Table~\ref{table:control_gen} indicate that both OFT and COFT achieve significantly stronger and more precise control than other approaches. We note that adapter-based methods (e.g., T2I-Adapter and ControlNet) tend to converge more slowly and produce inferior final outcomes. When compared to ControlNet, LoRA delivers better performance on the S2I task but underperforms on the C2I and L2F tasks. Overall, we observe that the performance limit of LoRA remains relatively low despite thorough hyperparameter tuning. In contrast, the performance of our OFT framework has not yet reached saturation, as we empirically observe continued improvement with an increasing number of trainable parameters. We highlight that our quantitative evaluation of controllable generation represents one of our key contributions, as it enables precise assessment of the control effectiveness of finetuned models (subject to the accuracy of the utilized off-the-shelf segmentation/detection models).

\textbf{Qualitative comparison}. We further conduct a qualitative comparison of OFT and COFT against leading state-of-the-art approaches, such as ControlNet, T2I-Adapter, and LoRA. The randomly generated images presented in Figure~\ref{fig:control_final} illustrate that OFT and COFT not only deliver high-fidelity and realistic image quality but also provide precise control. In the S2I task, it is evident that LoRA completely fails to produce images consistent with the input segmentation map, whereas ControlNet, OFT, and COFT effectively control the image generation. Compared to ControlNet, both OFT and COFT generate images with higher fidelity, richer details, and more plausible geometric structures while using significantly fewer model parameters. For the C2I task, OFT and COFT successfully hallucinate semantically similar images from rough Canny edges, whereas T2I-Adapter and LoRA exhibit much poorer performance. In the L2F task, our approach achieves the most precise pose control for generated faces, even under challenging facial poses. Across all three control tasks, we demonstrate that OFT and COFT qualitatively outperform the current state-of-the-art baselines, confirming the effectiveness of the OFT framework in controllable image generation. To offer a more thorough qualitative assessment, additional examples for all three control tasks are provided in Appendix~\ref{app:control_exp}. Furthermore, we showcase that OFT performs well on a broader set of control tasks, including dense pose to human body, sketch to image, and depth to image, detailed in Appendix~\ref{app:more_control_task}.

\section{Concluding Remarks and Open Problems}

Inspired by the insight that angular relationships among neurons play a key role in determining visual semantics, we introduce a straightforward yet powerful finetuning approach—orthogonal finetuning (OFT)—for managing text-to-image diffusion models. Our focus is on two specific text-to-image tasks: subject-driven generation and controllable generation. Compared with existing techniques, OFT offers enhanced controllability and greater stability during finetuning while requiring fewer finetuning parameters. Crucially, OFT does not add any inference time overhead, resulting in a model that is both effective and efficient to deploy.

OFT also raises several intriguing open questions. First, the orthogonality in OFT is ensured through Cayley parametrization, which necessitates a matrix inversion. This aspect somewhat restricts the scalability of OFT. Although we mitigate this issue by employing a block diagonal parametrization, finding a faster method to perform this matrix inversion in a differentiable manner remains an open challenge. Second, OFT shows promising potential for compositionality because the orthogonal matrices derived from multiple OFT finetuning tasks can be multiplied, preserving orthogonality. Investigating whether this group of orthogonal matrices retains the knowledge across all downstream tasks is a compelling research direction. Lastly, the parameter efficiency of OFT heavily relies on the block diagonal structure, which inevitably introduces biases and constrains flexibility. Developing approaches to enhance parameter efficiency more effectively and with fewer biases continues to be a significant open problem.

\end{document}